% -*- coding: utf-8 -*-

\chapter{Background and Related Work}

\section{Log Parsing and Template Mining}

Raw log messages are free-form text that combine natural language with runtime variable tokens
such as IP addresses, block identifiers, and numeric values. Before any pattern-based or
learning-based analysis can be applied, these messages must be transformed into structured event
templates by masking the variable portions. This task is known as log parsing, and it serves as a
prerequisite to virtually all downstream log analytics
pipelines~\cite{zhu2019tools}.

Several parsing approaches have been proposed with different tradeoffs between accuracy, throughput,
and adaptability to new log formats. Drain organizes discovered templates in a fixed-depth parse
tree, enabling efficient online matching and incremental template
updates~\cite{he2017drain}. Its fixed-depth design ensures that the parse tree remains compact
even when the number of templates grows, making it well-suited for continuous log streams. Spell
takes a streaming approach, using longest common subsequence matching to maintain message types as
new logs arrive without requiring global re-clustering~\cite{du2016spell}. Earlier work such as
IPLoM employs iterative partitioning, clustering messages by progressively splitting on token
positions~\cite{makanju2009iplom}. A systematic benchmark of parsing tools confirms that different
parsers exhibit substantially different accuracy and efficiency profiles on different log
corpora~\cite{zhu2019tools}, making parser selection a practically important preprocessing decision.
In our system, the Log Parser Agent performs a semantic variant of this task using an LLM, which
allows it to handle heterogeneous log formats without requiring a pre-defined template vocabulary.

\section{Root Cause Analysis: Traditional Approaches}

It is important to distinguish log-based RCA from the closely related but distinct task of anomaly
detection. Anomaly detection determines whether a sequence of log events or metrics deviates from
a learned or rule-defined model of normal behavior; it produces a binary or scored label
characterizing whether an execution window is anomalous~\cite{du2017deeplog,zhang2019logrobust}.
RCA, by contrast, is concerned with explaining \emph{why} an incident occurred: which component or
sub-system is the root cause, how the fault propagated, and what remediation is appropriate.
RCA takes anomaly detection as a possible input signal---it may reason over sequences of anomalous
events---but its objective is causal attribution rather than deviation
scoring~\cite{meng2020localizing,wu2020microrca}. This thesis is concerned exclusively with the
RCA task; anomaly detection methods are reviewed here only to the extent that they inform the log
representation strategies used by our parser.

Traditional automated RCA approaches can be broadly categorized into two families. The first family
employs causal graph methods: a dependency graph among services or metrics is constructed, and
statistical or graph-theoretic algorithms identify the component whose anomalous behavior is most
causally associated with the observed incident. CauseInfer builds a hierarchical causality graph
from service performance metrics to support automatic fault
localization~\cite{chen2014causeinfer}. MicroRCA uses a service call graph weighted by anomalous
performance indicators to localize root causes in microservice
systems~\cite{wu2020microrca}. CloudRanger applies Bayesian network learning to system metric
data to identify the most likely causal node in the dependency
graph~\cite{wang2018cloudranger}. The localization accuracy of these methods depends critically
on the quality of the dependency model; when the model is incomplete or stale, fault propagation
paths may not be correctly captured. Moreover, constructing and maintaining accurate causal graphs
requires ongoing engineering effort and deep system expertise, which limits practical deployment
in large, rapidly evolving cloud environments.

The second family applies supervised machine learning to log-derived representations. Early work
by Xu et al.\ demonstrates that principal component analysis over log message count matrices can
detect and localize problems in large-scale HDFS deployments, establishing logs as a viable signal
for automated system problem detection~\cite{xu2009consolelogs}. Subsequent work trains sequence
models over event template sequences to recognize fault signatures: DeepLog uses LSTM networks to
model normal execution sequences and flags out-of-distribution events as
anomalous~\cite{du2017deeplog}. While these methods can achieve strong detection rates on
controlled benchmarks, they require labeled training data from the target system, and their
performance degrades when log statement formatting changes across software versions---a pervasive
practical problem demonstrated empirically by LogRobust~\cite{zhang2019logrobust}. Furthermore,
detection-oriented classifiers produce a fault label or anomaly score but do not generate an
interpretable causal explanation, which is precisely what incident response workflows
require~\cite{chen2024rcacopilot}.

\section{Datasets and Benchmarks}

Public benchmarks have been essential for enabling reproducible evaluation and comparison across
log analysis approaches. LogHub aggregates log data from more than fifteen distinct software
systems---including Hadoop, HDFS, OpenStack, Apache, and Android---and has become the standard
reference for both anomaly detection and RCA research~\cite{he2020loghub}. The Hadoop1 and
HDFS\_v1 datasets used in our evaluation are drawn from LogHub. The HDFS dataset traces back to
foundational work demonstrating that console log analysis can scale to detecting problems across
tens of thousands of nodes~\cite{xu2009consolelogs}. For cloud management platforms, the CMCC
dataset provides industrial OpenStack failure traces in which failures propagate across interacting
services such as Nova, Neutron, and RabbitMQ; we use this dataset to evaluate cross-domain
generalization to real enterprise incidents~\cite{cui2025aetherlog}.

\section{Knowledge Graphs for Operational Reasoning}

Knowledge graphs represent entities, their attributes, and the relations between them in a
structured, queryable form. As a general technology, they support entity-centric retrieval,
multi-hop relational reasoning, and the integration of heterogeneous evidence
sources~\cite{hogan2021knowledge}. These properties are directly useful in AIOps: a KG that
stores historical incidents, the entities they involved, and their diagnosed root causes provides
a structured operational memory that can inform diagnosis of new incidents whose entities overlap
with past cases.

AetherLog is the most directly related prior work that combines KG retrieval with LLM-based
reasoning for log-based RCA~\cite{cui2025aetherlog}. By querying a KG of historical incidents and
injecting the retrieved context into the LLM prompt, AetherLog demonstrates that structured
historical grounding substantially reduces unsupported speculation and improves precision compared
to a purely parametric LLM baseline. Our work extends this direction by embedding KG retrieval
within a multi-agent debate framework, so that the retrieved context is evaluated critically
rather than accepted uncritically by a single reasoning pass.

\section{Retrieval-Augmented Generation}

Retrieval-augmented generation (RAG) augments a language model's parametric knowledge with
evidence retrieved from an external corpus at inference
time~\cite{lewis2021retrievalaugmentedgenerationknowledgeintensivenlp}. For knowledge-intensive
tasks where the relevant information may not have been adequately captured during pretraining, RAG
substantially improves factual accuracy by conditioning generation on retrieved documents. In the
RCA setting, retrieved evidence takes the form of incident reports, historical log summaries, or
structured KG subgraphs. The RAG paradigm motivates our KG Retrieval Agent, which supplies the
reasoning agents with entity-matched historical context. However, a limitation of standard RAG
pipelines is that retrieval quality is not verified: when retrieved context is irrelevant or
misleading, a single-agent model can still produce a confident but incorrect diagnosis. Roy et
al.\ explicitly identify this as a failure mode of LLM-based RCA agents, finding that
out-of-distribution retrieval can actively harm diagnostic
accuracy~\cite{roy2024llmagentRCA}. Our debate protocol addresses this by requiring the judge to
score evidence quality, penalizing hypotheses that cite weak or irrelevant retrieved facts.

\section{LLM-Based RCA: Current State of the Art}

The application of large language models to automated RCA has emerged as an active research
direction over the past two years, motivated by the observation that LLMs can interpret
heterogeneous log text and generate natural-language explanations without requiring
dataset-specific training or hand-crafted
rules~\cite{brown2020languagemodelsfewshotlearners,ouyang2022traininglanguagemodelsfollow}.

RCACopilot, deployed at Microsoft on Azure cloud incidents, represents the most prominent
production-scale LLM-based RCA system to date~\cite{chen2024rcacopilot}. It employs per-alert
incident handlers that collect diagnostic data from logs, traces, and metrics, and feeds this
aggregated context to an LLM for root cause category prediction. Evaluated on a year's worth of
production incidents, RCACopilot achieves up to 76.6\% RCA accuracy and demonstrates that LLMs
can match or exceed the accuracy of rule-based on-call systems when given sufficiently curated
diagnostic inputs. A key design choice is the separation of data aggregation (handled by
specialized handlers per alert type) from reasoning (handled by the LLM), which reduces the
context burden on the model. However, RCACopilot operates as a single-agent system with no
cross-validation of its prediction, and its reliance on specialized per-alert handlers limits
generalizability to new alert types without additional engineering.

Ahmed et al.\ conduct an early large-scale empirical study of GPT-3.x models for recommending
root causes and mitigation steps from Microsoft incident reports, finding that LLM outputs are
useful even without fine-tuning but that their quality degrades substantially for incidents whose
diagnostic signals do not appear in the model's training
distribution~\cite{ahmed2023rcarecommend}. This finding motivates the use of retrieval-based
grounding to supply in-context evidence that the model cannot rely on memorized knowledge to
provide.

Roy et al.\ study LLM-based ReAct agents equipped with retrieval tools for RCA on an
out-of-distribution dataset of production incidents~\cite{roy2024llmagentRCA}. Their evaluation
reveals a fundamental tension: agentic tool use can improve diagnostic coverage by retrieving
relevant evidence, but it also introduces new failure modes when the retrieval system returns
irrelevant results that the agent uncritically incorporates into its reasoning. This finding
directly informs our design: rather than relying on a single agent to both retrieve and reason,
we separate the retrieval function into a dedicated KG Retrieval Agent and subject the reasoning
outputs to competitive evaluation by the judge.

AetherLog specifically targets log-based RCA and integrates KG retrieval within the
LLM-based diagnostic pipeline~\cite{cui2025aetherlog}. Unlike the cloud incident systems of
RCACopilot and Ahmed et al., which consume pre-processed diagnostic bundles, AetherLog operates
on raw log sequences, making it the most architecturally comparable prior work to ours. Our
system extends AetherLog's core insight---that structured historical grounding improves LLM-based
log diagnosis---by adding the multi-agent debate and judging layers that AetherLog lacks.

\section{Foundations of Large Language Models}

Understanding why LLMs are effective for log-based reasoning requires reviewing the architectural
and training advances that underpin them.

\subsection{Transformer architectures}

The transformer architecture replaces recurrence with token-level self-attention, enabling parallel
computation over sequences and flexible modeling of long-range
dependencies~\cite{vaswani2023attentionneed}. This architectural property is directly relevant
to log analysis: failure signatures are often sparse, with the critical error token appearing many
lines after the triggering event, and self-attention provides a mechanism to associate these
distant tokens without the vanishing-gradient problems of recurrent models. BERT demonstrates that
deep bidirectional transformers pretrained on large corpora learn powerful contextual
representations applicable to a wide range of downstream
tasks~\cite{devlin2019bertpretrainingdeepbidirectional}, while RoBERTa establishes best practices
for pretraining data and hyperparameter choices that substantially improve representation
quality~\cite{liu2019robertarobustlyoptimizedbert}. Autoregressive decoder models such as
GPT-3 demonstrate that sufficiently large transformers can perform complex reasoning and
multi-step generation in a few-shot prompting setting, without any task-specific gradient
updates~\cite{brown2020languagemodelsfewshotlearners}. For RCA, this means that a capable
decoder LLM can generate structured hypotheses, evidence lists, and remediation narratives in
a single forward pass, conditioned solely on the prompt content.

\subsection{Scaling laws and emergent capabilities}

Empirical scaling laws establish that language model performance improves predictably as a
power-law function of model size, dataset size, and training
compute~\cite{kaplan2020scalinglawsneurallanguage}. Chinchilla scaling analysis further refines
this understanding, arguing that most existing LLMs were under-trained relative to their size and
that compute-optimal training requires substantially more data per parameter than earlier
recipes~\cite{hoffmann2022trainingcomputeoptimallargelanguage}. Large-scale training experiments
such as Gopher and PaLM characterize how emergent capabilities---including multi-step reasoning,
code generation, and in-context learning---appear at scale and improve non-monotonically with
model size~\cite{rae2022scalinglanguagemodelsmethods,chowdhery2022palmscalinglanguagemodeling}.
These findings have practical implications for our work: they explain why modern 7B-class open
models, which benefit from compute-efficient training recipes derived from scaling research, can
exhibit strong semantic generalization over heterogeneous operational text even in a zero-shot
setting.

\subsection{Instruction tuning and alignment}

Raw pretrained LLMs generate fluent text but do not reliably follow task instructions, adhere to
output format constraints, or avoid overconfident assertions. Instruction tuning and alignment
methods are designed to instill these behaviors. InstructGPT demonstrates that training with
reinforcement learning from human feedback (RLHF)---using PPO to optimize against a learned
reward model of human preferences---substantially improves instruction-following quality compared
to supervised fine-tuning
alone~\cite{ouyang2022traininglanguagemodelsfollow,schulman2017proximalpolicyoptimizationalgorithms}.
Complementary instruction-tuning approaches demonstrate that scale is not the only route to
zero-shot task generalization: finetuned language models evaluated on held-out task types show
strong zero-shot performance when instruction coverage in fine-tuning is
broad~\cite{wei2022finetunedlanguagemodelszeroshot}. Self-Instruct shows that self-generated
instructions can bootstrap alignment data at low cost~\cite{wang2023selfinstructaligninglanguagemodels},
while LIMA demonstrates that a small, carefully curated set of alignment examples can be
surprisingly effective~\cite{zhou2023limaalignment}. Direct preference optimization (DPO) offers
a simplified alternative to RLHF that avoids explicit reward model training by directly optimizing
the language model on preference
data~\cite{rafailov2024directpreferenceoptimizationlanguage}.

In the context of RCA, alignment is practically important in three respects. First, the system
must produce outputs in a strict JSON schema; misformatted outputs cannot be parsed and disrupt
the pipeline at scale. Second, the judge agent must produce calibrated scores rather than simply
asserting that one hypothesis is ``better.'' Third, reasoning agents must acknowledge uncertainty
rather than generating overconfident diagnoses when evidence is weak. Modern instruction-tuned
LLMs provide a strong foundation for these requirements, though reliability for operational tasks
still requires the additional architectural guardrails that our debate protocol provides.

\subsection{Parameter-efficient adaptation}

Even well-aligned foundation models may not have been exposed to the specific log terminology,
error patterns, and operational conventions of a target system during pretraining. Parameter-efficient
fine-tuning (PEFT) provides a practical mechanism for domain
adaptation~\cite{raffel2023exploringlimitstransferlearning}: LoRA inserts low-rank decomposition
matrices into transformer layers, adapting model behavior with a fraction of the parameters and
compute required for full fine-tuning~\cite{hu2021loralowrankadaptationlarge}. QLoRA further
reduces the resource footprint by combining quantization of the base model with LoRA adapters,
enabling fine-tuning of 7B-class models on commodity
hardware~\cite{dettmers2023qloraefficientfinetuningquantized}. While we do not perform
domain-specific fine-tuning in this thesis, these methods represent the most practical path for
future operational deployments where privacy constraints require on-premise adaptation on
organization-specific incident corpora.

\section{LLM Reasoning Strategies and Multi-Agent Verification}

Even a well-aligned, capable LLM can produce incorrect diagnoses when evidence is ambiguous or
incomplete. A growing body of work addresses this reliability problem through structured reasoning
strategies and multi-agent architectures.

Chain-of-thought prompting demonstrates that instructing a model to generate intermediate reasoning
steps before producing a final answer substantially improves performance on multi-step reasoning
tasks~\cite{wei2023chainofthoughtpromptingelicitsreasoning}. The key mechanism is that explicit
intermediate steps anchor the model's attention on the relevant sub-problems and reduce the chance
of shortcut reasoning that bypasses the evidentiary chain. Self-consistency extends this idea by
sampling multiple reasoning traces from the same model and selecting the most frequent final answer,
thereby reducing sensitivity to any single sampled
trajectory~\cite{wang2023selfconsistencyimproveschainthought}. Large LLMs are also zero-shot
reasoners~\cite{kojima2023largelanguagemodelszeroshot}, meaning that step-by-step reasoning can
be elicited simply by appending ``Let's think step by step'' to a prompt. ReAct interleaves
reasoning with tool invocation, enabling models to alternate between deliberation and action in
interactive environments~\cite{yao2023reactsynergizingreasoningacting}; Roy et al.\ apply this
approach to RCA specifically~\cite{roy2024llmagentRCA}.

Multi-agent debate takes a different approach to reliability: rather than improving the reasoning
of a single model, it introduces multiple agents that independently generate answers and then
critique each other's responses. Du et al.\ provide the first systematic empirical study of
multi-agent debate, demonstrating that having multiple LLM instances propose and refine answers
through iterative debate reduces factual errors and improves reasoning
quality~\cite{du2023improvingfactualityreasoninglanguage}. AutoGen formalizes multi-agent
conversation and tool use within a general framework that supports heterogeneous agent roles,
custom conversation topologies, and human-in-the-loop
interaction~\cite{wu2023autogen}. These frameworks establish the architectural vocabulary that
we draw on in designing our debate protocol, adapted to the specific requirements of log-based
RCA: hypothesis generation must be grounded in log evidence and historical KG context, the judge
must score against an operational rubric, and the protocol must terminate within a bounded number
of rounds to remain practical.

\section{Summary and Research Gap}

The literature reviewed in this chapter reveals a clear progression from template-based log
parsers through statistical RCA methods to LLM-based diagnostic systems, and from single-agent
LLM approaches to agent frameworks with tool use and retrieval augmentation. At the same time,
it reveals a gap that this thesis addresses.

Traditional causal-graph and supervised ML approaches to RCA require either a well-maintained
dependency model or labeled training data from the target system, neither of which is reliably
available in cross-domain or rapidly evolving cloud
environments~\cite{zhang2024cloudRCAsurvey,zhang2019logrobust}. LLM-based RCA systems such as
RCACopilot and AetherLog demonstrate that LLMs can diagnose incidents without task-specific
training, but they share the structural limitation of single-agent, single-pass generation: there
is no mechanism to verify whether the produced hypothesis is supported by the available evidence
or to systematically explore alternative explanations before committing to a
diagnosis~\cite{chen2024rcacopilot,roy2024llmagentRCA}. Knowledge graph grounding improves
factuality by supplying external evidence, but retrieval alone does not prevent the model from
ignoring or misinterpreting the retrieved context~\cite{cui2025aetherlog}.

Multi-agent debate has been shown to improve factual accuracy and reasoning quality in general
LLM settings~\cite{du2023improvingfactualityreasoninglanguage}, but this approach has not
previously been combined with knowledge graph grounding and applied to log-based RCA. This thesis
fills that gap by proposing a system in which multi-agent hypothesis generation, KG-based
evidence retrieval, and explicit judge-based evaluation are integrated into a single coherent
pipeline, and by evaluating this design empirically on three publicly available incident
datasets.
